\documentstyle[% 


righttag% 


,11pt% 


,amssymb% 


,russian%               remove in Eng. 


,izvru%                 Set izven% in Eng. 


]{amsart} 


 


% 


 


\newcommand{\tts}[1]{} 


 


%\hoffset=-15mm%                  For Eng. 


\setcounter{page}{3} 


\god{2000} 


\nomer{6~(457)} %                        Remove in Eng. 


%\nomer{Vol.\,44, No.\,6}%               Set in Eng. 


%\str{\pageref{begin}--\pageref{end}}%  Set in Eng. 


\udk{517.518.822} 


 


\author{\it Д.В.\,Алексеев} 


% NB:   ^^ remove in Eng. 


\title{Приближение функций нескольких переменных\\ 


с весом Чебышева--Эрмита} 


\thanks{Работа выполнена при финансовой поддержке Российского фонда 


фундаментальных исследований (грант \No\,97-01-00010).} 


 


\date{} 


%\pagestyle{myheadings}%         Set in Eng. 


\pagestyle{plain} %              Remove in Eng. 


\begin{document} 


%\thispagestyle{plain}%          Set in Eng. 


\maketitle 


\label{begin} 


 


 


В данной работе доказаны прямые теоремы для приближения 


функций многих переменных ``углом'' и ``прямоугольником'' 


из алгебраических полиномов и 


обратная теорема для приближения ``углом''. Соответствующие 


результаты для функций одной переменной были получены ранее в работах 


[1] и [2], а для приближения ``углом'' в метрике $L_2$ --- в [3]. 


 


Введем обозначения. Пусть 


$\rho(x)=\exp\big(-\frac{1}{2}\sum\limits^n_{i=1}x^2_i\big)$ --- 


весовая функция, $L_{p,\rho}(\Bbb R^N)$ --- множество 


измеримых по Лебегу функций $N$ переменных $f(x_1,\dotsc,x_N)$ с 


конечной нормой 


$$ 


\|f\|_{p,\rho}= 


\begin{cases} 


\displaystyle\bigg(\int_{\Bbb R^N}|f(x)\rho(x)|^p dx\bigg)^{1/p} 


&\text{при }\;1\le p<\infty;\\[3mm] 


\underset{x\in\Bbb R^N}{\text{ess\,sup}}{}\,|f(x)\rho(x)| 


&\text{при }\;p=\infty. 


\end{cases} 


$$ 


 


Определим оператор обобщенного сдвига по $i$-й переменной 


$$ 


T^i_h f(x_1,\dotsc,x_N)=\frac{1}{\sqrt{\pi}} 


\int^\infty_{-\infty}f(x_1,\dotsc,x_{i-1},u, 


x_{i+1},\dotsc,x_N)\exp(-y^2)d y, 


$$ 


где $u=e^{-h}(x_i+y\sqrt{e^{2h}-1})$. Введем дифференциальный оператор 


$D_i=\frac{1}{2}\frac{\partial^2}{\partial x^2_i}- 


x_i\frac{\partial}{\partial x_i}$. Этот оператор 


задан на классе $\Lambda_p(D_i)$ функций $f\in L_{p,\rho}(\Bbb R^n)$, 


для которых существует функция $g$, 


эквивалентная $f$ и такая, что при почти всех 


$(x_1,\dotsc,x_{i-1},x_{i+1},\dotsc,x_N)$ существует $\frac{\partial 


g}{\partial x_i}$, абсолютно 


непрерывная по $x_i$ на любом конечном отрезке и при этом 


$\frac{\partial^2 g}{\partial x^2_i}, x_i\frac{\partial g}{\partial 


x_i}\in L_{p,\rho}$. Обозначим 


$$ 


\Lambda_p(D^{r_i}_i)=\{f\in\Lambda_p(D_i)\mid 


D_i f\in\Lambda_p(D^{r_i-1}_i)\},\quad 


\Lambda_p(D^{r_1\dotsc r_k}_{i_1\dotsc i_k})= 


\{f\in\Lambda_p(D^{r_k}_{i_k})\mid 


D^{r_k}_{i_k}f\in\Lambda_p 


(D^{r_1\dotsc r_{k-1}}_{i_1\dotsc i_{k-1}})\}. 


$$ 


Определим обобщенный модуль гладкости 


$$ 


\Omega^{r_1\dotsc r_k}_{i_1\dotsc i_k} 


(f,\delta_1,\dotsc,\delta_k)_{p,\rho}= 


\sup_{|h_j|\le\delta_j,\ j=1\div k} 


\|(I-T^{i_1}_{h_1})^{r_1}\cdot\dotsc\cdot 


(I-T^{i_k}_{h_k})^{r_k}f\|_{p,\rho}, 


$$ 


где $I$ -- тождественный оператор. Пусть $K$-функционал 


$$ 


K^{r_1\dotsc r_k}_{i_1\dotsc i_k} 


(f,\delta_1,\dotsc,\delta_k)_{p,\rho}= 


\inf_{g_{l_1\dotsc l_j} 


\in\Lambda_p(D^{r_1\dotsc r_k}_{i_{l_1}\dotsc i_{l_j}})} 


\{\|f-\wh\Sigma g_{l_1\dotsc l_j}\|_{p,\rho}+ 


\wh\Sigma\delta^{r_{l_1}}_{l_1}\dotsb 


\delta^{r_{l_j}}_{l_j}\|D^{r_{l_1}}_{i_{l_1}}\dotsc 


d^{r_{l_j}}_{i_{l_j}}g_{l_1\dotsc l_j}\|_{p,\rho}\},


$$ 


где $\wh\Sigma$ обозначает $\sum\limits^k_{j=1}\sum\limits_{1\le 


l_1<l_2<\dotsb<l_j\le k}$. 


Пусть $A\le c B$, где $c$ --- постоянная, зависящая, быть может, от 


$p$, $N$, $k$, $r_i$, 


если эти параметры входят в $A$ или $B$, и не зависящая от других 


параметров, входящих в выражения $A$ и $B$. В 


этом случае будем использовать обозначение $A\preceq B$ и говорить, 


что $A$ не превосходит $B$ с 


константой. 


 


\begin{thm} 


Пусть $1\le p\le\infty$, $f\in L_{p,\rho}$, $1\le i_1<\dotsb<i_k\le 


N$, $\delta_j\in(0,1)$, 


тогда выполнены 


неравенства 


$$ 


K^{r_1\dotsc r_k}_{i_1\dotsc i_k}(f,\delta_1,\dotsc,\delta_k)_{p,\rho} 


\preceq\Omega^{r_1\dotsc r_k}_{i_1\dotsc i_k} 


(f,\delta_1,\dotsc,\delta_k)_{p,\rho}\preceq 


K^{r_1\dotsc r_k}_{i_1\dotsc i_k} 


(f,\delta_1,\dotsc,\delta_k)_{p,\rho}. 


$$ 


\end{thm} 


 


Пусть $\bold P^{i_1\dotsc i_k}_{n_1\dotsc 


n_k}=\Big\{P=\sum\limits^{n_1-1}_{l_1=0}\dotsb 


\sum\limits^{n_k-1}_{l_k=0}c_{l_1\dotsc l_k}x^{l_1}_{i_1}\dotsc 


x^{l_k}_{i_k}\Big\}$, 


где $c_{l_1\dotsc l_k}$ не зависят от $x_{i_1},\dotsc,x_{i_k}$ и 


принадлежат $L_{p,\rho}(\Bbb R^{N-k})$ как функции от остальных $N-k$ 


переменных. Будем обозначать 


$$ 


\bold A^{i_1\dotsc i_k}_{n_1\dotsc n_k}=\bigg\{ 


A=\sum^k_{l=1}P^{i_l}_{n_l},\ \;\text{где} 


\ \; P^{i_l}_{n_l}\in\bold P^{i_l}_{n_l}\bigg\}. 


$$ 


 


Величина $E^{i_1\dotsc i_k}_{n_1\dotsc n_k}(f)_{p,\rho}= 


\inf\limits_{P\in\bold P^{i_1\dotsc i_k}_{n_1\dotsc 


n_k}}\|f-P\|_{p,\rho}$ 


называется наилучшим приближением ``прямоугольником'', а величина 


$Y^{i_1\dotsc i_k}_{n_1\dotsc n_k}(f)_{p,\rho}= 


\inf\limits_{A\in\bold A^{i_1\dotsc i_k}_{n_1\dotsc 


n_k}}\|f-A\|_{p,\rho}$ --- 


наилучшим приближением ``углом'' по переменным 


$x_{i_1},\dotsc,x_{i_k}$. 


 


\begin{thm} 


Пусть $1\le i_1<\dotsb<i_k\le N$, $1\le p\le\infty$, $f\in 


L_{p,\rho}(\Bbb R^N)$. Тогда 


$$ 


E^{i_1\dotsc i_k}_{n_1\dotsc n_k}(f)_{p,\rho}\preceq 


\sum^k_{j=1}\Omega^{r_j}_{i_j}\bigg(f,\frac{1}{n_j} 


\bigg)_{p,\rho}. 


$$ 


\end{thm} 


 


\begin{thm} 


Пусть $1\le i_1<\dotsb<i_k\le N$, $1\le p\le\infty$, $f\in 


L_{p,\rho}(\Bbb R^N)$. Тогда 


$$ 


Y^{i_1\dotsc i_k}_{n_1\dotsc n_k}(f)_{p,\rho}\preceq 


\Omega^{r_1\dotsc r_k}_{i_1\dotsc 


i_k}\bigg(f,\frac{1}{n_1},\dotsc,\frac{1}{n_k} \bigg)_{p,\rho}. 


$$ 


\end{thm} 


 


\begin{thm} 


Пусть $1\le i_1<\dotsb<i_k\le N$, $1\le p\le\infty$, $f\in 


L_{p,\rho}(\Bbb R^N)$. Тогда 


$$ 


\Omega^{r_1\dotsc r_k}_{i_1\dotsc 


i_k}\bigg(f,\frac{1}{n_1},\dotsc,\frac{1}{n_k} \bigg)_{p,\rho} 


\preceq\frac{1}{n^{r_1}_1\dotsc n^{r_k}_k} 


\sum^{n_1}_{s_1=1}\dotsb\sum^{n_k}_{s_k=1} 


s^{r_1-1}_1\dotsc s^{r_k-1}_k 


Y^{i_1\dotsc i_k}_{s_1\dotsc s_k}(f)_{p,\rho}. 


$$ 


\end{thm} 


 


{\bf Доказательство теоремы 1.} Для доказательства потребуются 


следующие леммы. 


 


\begin{lem} 


Пусть $h\in(0,1)$. Тогда $T^i_h$ --- ограниченный оператор, т.\,е. 


$\|T^i_h f\|_{p,\rho}\preceq\|f\|_{p,\rho}$ 


при $f\in L_{p,\rho}$. 


\end{lem} 


 


\begin{lem} 


Для любой $f\in L_{p,\rho}$ существует $g_i\in\Lambda_p(D_i)$ такая, 


что $D_i g_i=f-c_i$, где 


$c_i=c_i(x_1,\dotsc,x_{i-1},x_{i+1},\dotsc,x_N)\in L_{p,\rho}(\Bbb 


R^{N-1})$. 


\end{lem} 


 


\begin{lem} 


Пусть $g_i\in\Lambda_p(D_i)$. Тогда 


$$ 


(T^i_h-I)g_i=\int^h_0 T^i_\varphi D_i g_i d\varphi. 


$$ 


\end{lem} 


 


\begin{lem} 


При $h_1,h_2>0$, $h_1+h_2<1$ справедливо групповое свойство 


$$ 


T^i_{h_1}T^i_{h_2}=T^i_{h_1+h_2}. 


$$ 


\end{lem} 


 


\begin{lem} 


Пусть $f\in \Lambda_p(D^r_i)$, $h\in (0,1)$. Тогда 


$$ 


(T^i_h-I)^r f=h^r\int^r_0\Pi_k(s) T^i_{h s}D^r_i f\,d s, 


$$ 


где $\Pi_r(s)=\sum\limits^r_{u=0}(-1)^{r+1}C^u_r 


\frac{(u-s)^{r-1}_+}{(r-1)!}$ --- ядро Пеано, $(a)_+=\max\{a,0\}$. 


\end{lem} 


 


\begin{lem} 


Пусть $f\in L_{p,\rho}$, $h\in(0,1)$. Тогда функция 


$g_r=\int\limits^r_0\Pi_r(s)T^i_{h s}f\,d s$ принадлежит 


$\Lambda_p(D^r_i)$, 


и при этом 


$$ 


D^r_i g_r=\frac{1}{h^k}(T^i_h-I)^k f. 


$$ 


\end{lem} 


 


\begin{lem} 


Пусть $1\le i_1<\dotsb<i_k\le N$. Тогда для любой 


$f\in\Lambda_p(D^{r_1\dotsc r_k}_{i_1\dotsc i_k})$ выполнено 


неравенство 


$$ 


\Omega^{r_1\dotsc r_k}_{i_1\dotsc i_k}(f;\delta)_{p,\rho}\preceq 


\delta^{r_1}_{i_1}\dotsc \delta^{r_k}_{i_k} 


\|D^{r_1}_{_1}\dotsc D^{r_k}_{i_l}f\|_{p,\rho}. 


$$ 


\end{lem} 


 


Леммы 1--7 позволяют доказать теорему 1. 


Для наглядности теорема будет доказана в случае 


$N=k=2$. Докажем левое неравенство. Из обобщенного неравенства 


треугольника для модулей 


гладкости имеем 


$$ 


\Omega^{r_1r_2}_{p,\rho}(f;\delta_1\delta_2)\preceq 


\Omega^{r_1r_2}_{p,\rho}(f-g_1-g_2-g_{12};\delta_1\delta_2)+ 


\Omega^{r_1r_2}_{p,\rho}(g_1;\delta_1\delta_2)+ 


\Omega^{r_1r_2}_{p,\rho}(g_2;\delta_1\delta_2)+ 


\Omega^{r_1r_2}_{p,\rho}(g_{12};\delta_1\delta_2) 


$$ 


для произвольных 


\begin{equation} 


g_1\in\Lambda_p(D^{r_1}_1),\ g_2\in\Lambda_p(D^{r_2}_2), 


\ g_{12}\in\Lambda_p(D^{r_1r_2}_{12}). 


\end{equation} 


Оценим первое слагаемое по лемме 1 как $\|f-g_1-g_2-g_{12}\|_{p,\rho}$, 


а остальные --- по лемме 7. 


Получим 


$$ 


\Omega^{r_1r_2}_{p,\rho}(f;\delta_1\delta_2)\preceq 


\|f-g_1-g_2-g_{12}\|_{p,\rho}+\delta^{r_1}_1 


\|D^{r_1}_1g_1\|_{p,\rho}+ 


\delta^{r_2}_2\|D^{r_2}_2g_2\|_{p,\rho}+ 


\delta^{r_1}_1\delta^{r_2}_2\|D^{r_1}_1 D^{r_2}_2\|_{p,\rho}. 


$$ 


Взяв точную нижнюю грань от правой части по всем $g_1$, $g_2$, 


$g_{12}$, 


удовлетворяющим (1), получим 


определение $K$-функционала, что и требовалось доказать. 


 


Докажем правое неравенство. Рассмотрим функции вида 


\begin{align*} 


g_1&=\int^{r_1}_0\int^{r_2}_0\Pi_{r_1}(t_1)\Pi_{r_2}(t_2) 


\big(I-(I-T^1_{\frac{\delta_1t_1}{r_1}})^{r_1}\big) 


(I-T^2_{\frac{\delta_2t_2}{r_2}})^{r_2}f\,dt_1dt_2,\\ 


% 


g_2&=\int^{r_1}_0\int^{r_2}_0\Pi_{r_1}(t_1)\Pi_{r_2}(t_2) 


\big(I-(I-T^2_{\frac{\delta_2t_2}{r_2}})^{r_2}\big) 


(I-T^1_{\frac{\delta_1t_1}{r_1}})^{r_1}f\,dt_1dt_2,\\ 


% 


g_{12}&=\int^{r_1}_0\int^{r_2}_0\Pi_{r_1}(t_1)\Pi_{r_2}(t_2) 


\big(I-(I-T^1_{\frac{\delta_1t_1}{r_1}})^{r_1}\big) 


(I-(I-T^2_{\frac{\delta_2t_2}{r_2}})^{r_2})f\,dt_1dt_2,


\end{align*} 


где $\Pi_r(t)$ --- ядро Пеано. Используя тот факт, 


что $\int\limits^r_0\Pi_r(t)dt=1$ ([4], c.\,15--16), получим 


$$ 


f-g_1-g_2-g_{12}=\int^{r_1}_0\int^{r_2}_0\Pi_{r_1}(t_1) 


\Pi_{r_2}(t_2)\big(I-T^1_{\frac{\delta_1t_1}{r_1}}\big)^{r_1} 


\big(I-T^2_{\frac{\delta_2t_2}{r_2}}\big)^{r_2}f\,dt_1dt_2. 


$$ 


Отсюда, применив обобщенное неравенство треугольника и пользуясь 


доказанной в [4] неотрицательностью ядра Пеано, получим 


$$ 


\|f-g_1-g_2-g_{12}\|_{p,\rho}\le 


\Omega^{r_1r_2}_{p,\rho}(f;\delta_1\delta_2). 


$$ 


 


Теперь оценим $\|D^{r_1}_1g_1\|_{p,\rho}$ (по лемме 6\; 


$g_1\in\Lambda_p(D^{r_1}_1)$). Применив лемму 6 с $h=s_1\delta_1/r_1$, 


получим 


$$ 


D^{r_1}_1g_1=\frac{1}{\delta^{r_1}_1}\sum^{r_1}_{s_1=1} 


(-1)^{r_1+s_1+1}C^{s_1}_{r_1}\bigg(\frac{r_1}{s_1} \bigg)^{r_1}


\int^{r_2}_0\big(I-T^1_{\frac{\delta_1s_1}{r_1}}\big)^{r_1} 


\big(I-T^2_{\frac{\delta_2t_2}{r_2}}\big)^{r_2}f\,dt_2. 


$$ 


Оценивая норму этого выражения по обобщенному неравенству 


Минковского и домножая на $\delta^{r_1}_1$, 


установим, что 


$$ 


\delta^{r_1}_1\|D^{r_1}_1g_1\|_{p,\rho}\preceq 


\Omega^{r_1r_2}_{p,\rho}(f;\delta_1\delta_2). 


$$ 


Проводя аналогичные выкладки для $g_2$, получим 


$$ 


\delta^{r_2}_2\|D^{r_2}_2g_2\|_{p,\rho}\preceq 


\Omega^{r_1r_2}_{p,\rho}(f;\delta_1\delta_2). 


$$ 


 


Оценим теперь $\|D^{r_1}_1D^{r_2}_2g_{12}\|_{p,\rho}$. Применяя дважды 


лемму 6 (сначала с $h=s_1\delta_1/r_1$, а потом с 


$h=s_2\delta_2/r_2$), имеем 


$$ 


D^{r_1}_1D^{r_2}_2g_{12}=\frac{1}{\delta^{r_1}_1\delta^{r_2}_2} 


\sum^{r_1}_{s_1=1}\sum^{r_2}_{s_2=1} 


a(r_1,r_2,s_1,s_2)\big(I-T^1_{\frac{\delta_1s_1}{r_1}}\big)^{r_1} 


\big(I-T^2_{\frac{\delta_2s_2}{r_2}}\big)^{r_2}f, 


$$ 


где $a(r_1,r_2,s_1,s_2)$ --- коэффициенты, не зависящие от $f$ и 


$\delta_i$. Умножая на $\delta^{r_1}_1\delta^{r_2}_2$ и оценивая 


норму по обобщенному неравенству треугольника, получаем 


$$ 


\delta^{r_1}_1\delta^{r_2}_2\|D^{r_1}_1D^{r_2}_2g_{12}\|_{p,\rho} 


\preceq\Omega^{r_1r_2}_{p,\rho}(f;\delta_1\delta_2). 


$$ 


 


Таким образом, 


$$ 


\|f-g_1-g_2-g_{12}\|_{p,\rho}+\delta^{r_1}_1 


\|D^{r_1}_1g_1\|_{p,\rho}+ 


\delta^{r_2}_2\|D^{r_2}_2g_2\|_{p,\rho}+ 


\delta^{r_1}_1\delta^{r_2}_2\|D^{r_1}_1D^{r_2}_2g_{12}\|_{p,\rho} 


\preceq\Omega^{r_1r_2}_{p,\rho}(f;\delta_1\delta_2). 


$$ 


Если эта оценка верна для конкретных $g_1$, $g_2$, $g_{12}$, 


то для нижней грани выражения она тем 


более выполнена. Следовательно, из определения $K$-функционала 


вытекает правое неравенство.$\quad\square$ 


 


Для доказательства теорем 2, 3 потребуются леммы 8--13. 


 


\begin{lem} 


Пусть $1\le i_1<\dotsb<i_k\le N$. Если 


$P\in\bold P^{i_1\dotsc i_k}_{n_1\dotsc n_k}$, то выполнено неравенство 


$$ 


\|D^{r_1}_{i_1}\dotsc D^{r_k}_{i_k}P\|_{p,\rho}\preceq 


n^{r_1}_1\dotsc n^{r_k}_k\|P\|_{p,\rho}. 


$$ 


\end{lem} 


 


Пусть $V^i_m f$ --- средние Валле-Пуссена ряда Фурье--Эрмита функции 


$f$ по переменной $x_i$ и 


$\pi^i_n f=V^i_{[n/2]}f$. 


 


\begin{lem} 


Пусть $f\in L_{p,\rho}$. Тогда 


$\|\pi^i_n f\|_{p,\rho}\preceq\|f\|_{p,\rho}$.


\end{lem} 


 


\begin{lem} 


Если $f\in L_{p,\rho}$, то $\pi^i_n f\in\bold P^i_n$. 


\end{lem} 


 


\begin{lem} 


Если $f\in\Lambda_p(D_i)$, то 


$$ 


\|\pi^i_n f-f\|_{p,\rho}\preceq\frac{1}{n} 


\|D_i f\|_{p,\rho}. 


$$ 


\end{lem} 


 


\begin{lem} 


Пусть $1\le i_1<\dotsb<i_k\le N$ и пусть 


$P\in\bold P^{n_1\dotsc n_k}_{i_1\dotsc i_k}$, т.\,е. 


$$ 


P=\sum^{n_1-1}_{s_1=0}\dotsc\sum^{n_k-1}_{s_k=0} 


c_{s_1\dotsc s_k}x^{s_1}_{i_1}\dotsc x^{s_k}_{i_k},\ \; 


c_{s_1\dotsc s_k}\in L_{p,\rho}(\Bbb R^{N-k}). 


$$ 


Тогда $\Omega^{r_1\dotsc r_k}_{i_1\dotsc i_k}(P;\delta)\preceq 


(n_1\delta_1)^{r_1}\dotsb(n_k\delta_k)^{r_k}\|P\|_{p,\rho}$. 


\end{lem} 


 


\begin{lem} 


Для любой $f\in L_{p,\rho}$ выполнено неравенство 


$$ 


\|(I-\pi^1_{2n})(I-\pi^2_{2m})f\|_{p,\rho}\preceq 


Y_{n m}(f)_{p,\rho}. 


$$ 


\end{lem} 


 


Теоремы 2, 3 будут доказаны в случае $N=k=2$. 





{\bf Доказательство теоремы 2.}  


Рассмотрим 


$P_{n_1n_2}(f)=(I-(I-\pi^1_{n_1})^{r_1})(I-(I-\pi^2_{n_2})^{r_2})$. 


По лемме 10\; $P_{n_1n_2}\in\bold P^{12}_{n_1n_2}$. Следовательно, 


$E^{12}_{n_1n_2}(f)_{p,\rho}\le\|f-P_{n_1n_2}\|_{p,\rho}$. По 


неравенству треугольника 


$$ 


\|f-P_{n_1n_2}\|_{p,\rho}\le\|(I-\pi^1_{n_1})^{r_1}f\|_{p,\rho}+ 


\|(I-(I-\pi^2_{n_2})^{r_2})f\|_{p,\rho}+ 


\|(I-(I-\pi^1_{n_1})^{r_1})(I-(I-\pi^2_{n_2})^{r_2} 


)f\|_{p,\rho}, 


$$ 


а это по лемме 9 не превосходит с константой 


величины $\|(I-\pi^1_{n_1})^{r_1}\|_{p,\rho}+ 


\|(I-\pi^2_{n_2})^{r_2}f\|_{p,\rho}$. 


Значит,


\begin{equation} 


E_{n_1n_2}(f)\preceq\|(I-\pi^1_{n_1})^{r_1}f\|_{p,\rho}+ 


\|(I-\pi^2_{n_2})^{r_2}f\|_{p,\rho}. 


\end{equation} 


Рассмотрим произвольную $g_1\in\Lambda_p(D^{r_1}_1)$. Тогда 


$$ 


\|(I-\pi^1_{n_1})^{r_1}f\|_{p,\rho}\le 


\|(I-\pi^1_{n_1})^{r_1}(f-g_1)\|_{p,\rho}+ 


\|(I-\pi^1_{n_1})^{r_1}g_1\|_{p,\rho}. 


$$ 


Оценивая первое слагаемое по лемме 9, а второе --- по лемме 11, получим 


$$ 


\|(I-\pi^1_{n_1})^{r_1}f\|_{p,\rho}\preceq 


\|f-g_1\|_{p,\rho}+\frac{1}{n^{r_1}_1} 


\|D^{r_1}_1g_1\|_{p,\rho}. 


$$ 


Взяв инфимум по всем $g_1\in\Lambda_p(D^{r_1}_1)$, получим 


$$ 


\|(I-\pi^1_{n_1})^{r_1}f\|_{p,\rho}\preceq 


K^{r_1}_1(f,1/n)_{p,\rho}. 


$$ 


Применяя теорему 1, имеем 


$$ 


\|(I-\pi^1_{n_1})^{r_1}f\|_{p,\rho}\preceq 


\Omega^{r_1}_1(f,1/n_1)_{p,\rho}. 


$$ 


Аналогично доказывается, что 


$\|(I-\pi^2_{n_2})^{r_2}f\|_{p,\rho}\preceq 


\Omega^{r_2}_2(f,1/n_2)_{p,\rho}$. 


Подстановкой этих оценок в (2) и завершается доказательство.$\quad\square$ 


 


{\bf Доказательство теоремы 3.} 


Рассмотрим $A_{n_1n_2}=f-L f$, где $L$ --- 


оператор вида $L f=(I-\pi^1_{n_1})^{r_1}(I-\pi^2_{n_2})^{r_2}$. 


Согласно лемме 10\; $A_{n_1n_2}\in\bold A^{12}_{n_1n_2}$. Следовательно, 


\begin{equation} 


Y^{12}_{n_1n_2}(f)_{p,\rho}\le 


\|f-A_{n_1n_2}\|_{p,\rho}=\|L f\|_{p,\rho}. 


\end{equation} 


Выберем произвольные $g_1\in\Lambda_p(D^{r_1}_1)$, 


$g_2\in\Lambda_p(D^{r_2}_2)$, $g_{12}\in\Lambda_p(D^{r_1r_2}_{12})$. Из 


неравенства треугольника имеем 


$$ 


\|L f\|_{p,\rho}\le\|L(f-g_1-g_2-g_{12})\|_{p,\rho} 


+\|L g_1\|_{p,\rho}+\|L g_2\|_{p,\rho}+ 


\|L g_{12}\|_{p,\rho}. 


$$ 


Оценивая слагаемые в правой части по леммам 9 и 11, получим 


$$ 


\|L f \|_{p,\rho}\preceq\|f-g_1-g_2-g_{12}\|_{p,\rho}+ 


\frac{1}{n^{r_1}_1}\|D^{r_1}_1g_1\|_{p,\rho}+ 


\frac{1}{n^{r_2}_2}\|D^{r_2}_2g_2\|_{p,\rho}+ 


\frac{1}{n^{r_1}_1n^{r_2}_2}\|D^{r_1}_1D^{r_2}_2g_{12}\|_{p,\rho}. 


$$ 


Взяв нижнюю грань правой части по всем $g_1\in\Lambda_p(D^{r_1}_1)$, 


$g_2\in\Lambda_p(D^{r_2}_2)$, $g_{12}\in\Lambda_p(D^{r_1r_2}_{12})$, 


получим 


$$ 


\|L f\|_{p,\rho}\preceq K^{r_1r_2}_{12} 


\bigg(f,\frac{1}{n_1},\frac{1}{n_2} \bigg). 


$$ 


По теореме 1 


$$ 


\|L f\|_{p,\rho}\preceq \Omega^{r_1r_2}_{12} 


\bigg(f,\frac{1}{n_1},\frac{1}{n_2} \bigg). 


$$ 


Подставив эту оценку в (3), завершим доказательство.$\quad\square$ 


 


{\bf Доказательство теоремы 4.} Рассмотрим выражение 


$$ 


\{I-\pi^1_1\}\{I-\pi^2_1\}f=\bigg\{I-\pi^1_{2^L}+ 


\sum^{L-1}_{l=0}(\pi^1_{2^{l+1}}-\pi^1_{2^l})\bigg\} 


\bigg\{I-\pi^2_{2^M}+\sum^{M-1}_{m=0} 


(\pi^2_{2^{m+1}}-\pi^2_{2^m})\bigg\}f. 


$$ 


Раскрывая фигурные скобки, получим 


\begin{multline} 


f-\pi^1_1-\pi^2_1+\pi^1_1\pi^2_1f=\sum^{L-1}_{l=0} 


\sum^{M-1}_{m=0}(\pi^1_{2^{l+1}}-\pi^1_{2^l}) 


(\pi^2_{2^{m+1}}-\pi^2_{2^m})f+\\ 


% 


+(I-\pi^2_{2^M})\sum^{L-1}_{l=0} 


(\pi^1_{2^{l+1}}-\pi^1_{2^l})f+(I-\pi^1_{2^L}) 


\sum^{M-1}_{m=0}(\pi^2_{2^{m+1}}-\pi^2_{2^m})f+ 


(I-\pi^1_{2^L})(I-\pi^2_{2^M})f. 


\end{multline} 


Из свойств оператора $\pi^j_{n_j}$ следует, что $\pi^1_1f$ зависит 


только от $x$, $\pi^2_1f$ --- только от $y$, а $\pi^1_1\pi^2_1f$ 


постоянна почти всюду. Следовательно, 


они не влияют на величину обобщенного модуля гладкости. 


Таким образом, для того чтобы оценить величину 


$\Omega^{r_1r_2}_{p,\rho}(f;\delta_1\delta_2)$, достаточно оценить 


модули 


гладкости слагаемых, стоящих в правой части (4). 


 


Первое слагаемое есть сумма выражений вида 


$\{\pi^1_{2^{l+1}}-\pi^1_{2^l}\}\{\pi^2_{2^{m+1}}-\pi^2_{2^m}\}f= 


\{(I-\pi^1_{2^l})-(I-\pi^1_{2^{l+1}})\}\{(I-\pi^2_{2^m})- 


(I-\pi^2_{2^{m+1}})\}f$. 


Раскрывая фигурные скобки и применяя лемму 13, нетрудно 


получить оценку 


$$ 


\|(\pi^1_{2^{l+1}}-\pi^1_{2^l})(\pi^2_{2^{m+1}}-\pi^2_{2^m})f 


\|_{p,\rho}\preceq Y_{2^{l-1}2^{m-1}}(f)_{p,\rho}. 


$$ 


Выражение, стоящее под знаком нормы, 


есть полином степени $2^{l+1}-1$ по $x$ и $2^{m+1}-1$ по $y$. 


Следовательно, по лемме 12 можно оценить выражение 


\begin{equation} 


\Omega^{r_1r_2}_{p,\rho}((\pi^1_{2^{l+1}}-\pi^1_{2^l}) 


(\pi^2_{2^{m+1}}-\pi^2_{2^m})f;\delta_1\delta_2)\preceq 


\delta^{r_1}_1\delta^{r_2}_2 2^{l r_1+m r_2} 


Y_{2^{l-1}2^{m-1}}(f)_{p,\rho}. 


\end{equation} 


Второе слагаемое в правой части (4) есть сумма выражений вида 


$$ 


(I-\pi^2_{2^M})(\pi^1_{2^{l+1}}-\pi^1_{2^l})f= 


(I-\pi^2_{2^M})(I-\pi^1_{2^l})f- 


(I-\pi^2_{2^M})(I-\pi^1_{2^{l+1}})f. 


$$ 


Применив к ним лемму 12, получим 


\begin{equation} 


\Omega^{r_1r_2}_{p,\rho}((I-\pi^2_{2^M})(\pi^1_{2^{l+1}}- 


\pi^1_{2^l})f;\delta_1\delta_2)\preceq 


\delta^{r_1}_1\delta^{r_2}_2 2^{l r_1+M r_2} 


Y_{2^{l-1}2^{M-1}}(f)_{p,\rho}. 


\end{equation} 


Аналогично оцениваются компоненты третьего слагаемого 


\begin{equation} 


\Omega^{r_1r_2}_{p,\rho}((I-\pi^1_{2^L}) 


(\pi^1_{2^{m+1}}-\pi^2_{2^m})f;\delta_1\delta_2)\preceq 


\delta^{r_1}_1\delta^{r_2}_2 2^{N r_1+m r_2} 


Y_{2^{L-1}2^{m-1}}(f)_{p,\rho},


\end{equation} 


а четвертое слагаемое просто оценим величиной 


\begin{equation} 


\Omega^{r_1r_2}_{p,\rho}((I-\pi^1_{2^l}) 


(I-\pi^2_{2^M})f;\delta_1\delta_2)\preceq 


Y_{2^{L-1}2^{M-1}}(f)_{p,\rho}. 


\end{equation} 


Взяв обобщенный модуль гладкости от 


обеих частей (4) и подставив оценки (5)--(8), получим 


\begin{multline} 


\Omega^{r_1r_2}_{p,\rho}(f;\delta_1\delta_2)\preceq 


\delta^{r_1}_1\delta^{r_2}_2\bigg\{\sum^{L-1}_{l=0} 


\sum^{M-1}_{m=0}2^{r_1l+r_2m}Y_{2^{l-1}2^{m-1}}(f)_{p,\rho}+\\ 


% 


+\sum^{N-1}_{l=0}2^{r_1l+r_2M}Y_{2^{l-1}2^{M-1}}(f)_{p,\rho}+ 


\sum^{M-1}_{m=0}2^{r_1L+r_2m}Y_{2^{L-1}2^{m-1}} 


(f)_{p,\rho}\bigg\}+Y_{2^{L-1}2^{M-1}}(f)_{p,\rho}. 


\end{multline} 


Подставив $\delta_1=2^{-L}$, $\delta_2=2^{-M}$ в (9), получим 


\begin{multline} 


\Omega^{r_1r_2}_{p,\rho}(f;\delta_1\delta_2)\preceq 


2^{-(r_1l+r_2M)}\bigg\{\sum^{L-1}_{l=0}\sum^{M-1}_{m=0} 


2^{l(r_1-1)}2^{m(r_2-1)}2^{l+m}Y_{2^{l-1}2^{m-1}}(f)_{p,\rho}+\\ 


% 


+\sum^{L-1}_{l=0}2^{(r_1-1)l+r_2M}2^l 


Y_{2^{l-1}2^{M-1}}(f)_{p,\rho}+ \\ 


% 


+\sum^{M-1}_{m=0}2^{r_1L+(r_2-1)m}2^m 


Y_{2^{L-1}2^{m-1}}(f)_{p,\rho}+ 


2^{r_1N+r_2M}Y_{2^{L-1}2^{M-1}}(f)_{p,\rho}\bigg\}. 


\end{multline} 


Очевидно, $\sum\limits^{S-1}_{s=0}2^{r s}\ge C_r2^{r S}$. 


Следовательно, заменив в оценке (10) $2^{r_1L}$ на 


$\sum\limits^{L-1}_{l=0}2^{r_1l}$ и $2^{r_2M}$ на 


$\sum\limits^{M-1}_{m=0}2^{r_2m}$, получим 


\begin{multline*} 


\Omega^{r_1r_2}_{p,\rho}(f;\delta_1\delta_2)\preceq \\ 


\preceq 2^{-(r_1L+r_2M)}\sum^{L-1}_{l=0} 


\sum^{M-1}_{m=0}2^{r_1l+r_2m} 


[Y_{2^{l-1}2^{m-1}}(f)_{p,\rho}+ 


Y_{2^{l-1}2^{M-1}}(f)_{p,\rho}+ 


Y_{2^{L-1}2^{m-1}}(f)_{p,\rho}+ 


Y_{2^{L-1}2^{M-1}}(f)_{p,\rho}]. 


\end{multline*} 


Так как $Y_{m n}$ есть невозрастающая функция по аргументам $m$ и $n$, 


то выражение, стоящее в 


квадратных скобках, не превосходит $4Y_{2^{l-1}2^{m-1}}(f)_{p,\rho}$. 


Следовательно, 


\begin{equation} 


\Omega^{r_1r_1}_{p,\rho}(f;\delta_1\delta_2)\preceq 


2^{-(r_1L+r_2M)}\sum^{L-1}_{l=0}\sum^{M-1}_{m=0} 


2^{r_1l+r_2m}Y_{2^{l-1}2^{m-1}}(f)_{p,\rho}. 


\end{equation} 


Аналогично можно получить оценку 


$$


2^{l+m}Y_{2^{l-1}2^{m-1}}(f)_{p,\rho}\preceq 


\sum^{2^{l-1}}_{s_1=2^{l-2}+1}\sum^{2^{m-1}}_{s_2=2^{m-2}+1} 


Y_{s_1s_2}(f)_{p,\rho}. 


$$


Подставив эти величины в (11), получим 


$$ 


\Omega^{r_1r_2}_{12}\bigg(f;\frac{1}{R_1},\frac{1}{R_2} \bigg) 


\preceq\frac{1}{R^{r_1}_1R^{r_2}_2}\sum^{R_1-1}_{s_1=1} 


\sum^{R_2-1}_{s_2=1}s^{r_1-1}_1s^{r_2-1}_2 Y_{s_1s_2}


(f)_{p,\rho}, 


$$ 


где $R_1=2^L$, $r_2=2^M$. Для того чтобы получить оценку для 


произвольных $R_1$, $R_2$, достаточно 


выбрать $M$ и $L$ так, что $2^L\le R_1<2^{L+1}$ и


$2^M\le R_2<2^{M+1}.\quad\square$ 


 


\begin{cor} 


Пусть $f\in L_{p,\rho}(R^N)$, $1\le i_1<\dotsb<i_k\le N$, 


$\delta_1,\dotsc,\delta_k\in(0,1)$, $\alpha_j<r_j$. 


Тогда 


$$ 


\theta^{r_1}_{i_1}(f,\delta_1)_{p,\rho}\preceq 


\delta^{\alpha_1}_1,\dotsc,\Omega^{r_k}_{i_k}(f,\delta_k 


)_{p,\rho}\preceq\delta^{\alpha_k}_k 


$$ 


в том и только том случае, когда 


$$ 


E^{i_1\dotsc i_k}_{n_1\dotsc n_k}(f)_{p,\rho}\preceq 


\sum^k_{j=1}\frac{1}{n^{\alpha_j}_j}. 


$$ 


\end{cor} 


 


\begin{cor} 


Пусть $f\in L_{p,\rho}(R^N)$, $1\le i_1<\dotsb<i_k\le N$, 


$\delta_1,\dotsc,\delta_k\in(0,1)$, $\alpha_j<r_j$. 


Тогда 


$$ 


\Omega^{r_1\dotsc r_k}_{i_1\dotsc i_k} 


(f,\delta_1,\dotsc,\delta_k)_{p,\rho}\preceq 


\delta^{\alpha_1}_1\dotsc\delta^{\alpha_k}_k 


$$ 


в том и только том случае, когда 


$$ 


Y^{i_1\dotsc i_k}_{n_1\dotsc n_k}(f)_{p,\rho}\preceq 


\frac{1}{n^{\alpha_1}_1\cdot\dotsc\cdot n^{\alpha_k}_k}. 


$$ 


\end{cor} 


 


{\bf Доказательство следствия 1.} Если 


$\Omega^{r_j}_{i_j}(f)_{p,\rho}\preceq\delta^{\alpha_j}_j$, то 


$E^{i_1\dotsc i_k}_{n_1\dotsc n_k}(f)_{p,\rho}\preceq 


\sum\limits^k_{j=1}n^{-\alpha_j}_j$ 


по теореме~2. Если же 


$E^{i_1\dotsc i_k}_{n_1\dotsc n_k}(f)_{p,\rho}\preceq 


\sum\limits^k_{j=1}n^{-\alpha_j}_j$, 


то, устремляя $n_j$ при всех $j\ne j_0$ к бесконечности, имеем 


$E^{r_{j_0}}_{n_{j_0}}(f)_{p,\rho}\preceq n^{-\alpha_{j_0}}_{j_0}$.


По определению


$E^{r_{j_0}}_{n_{j_0}}(f)_{p,\rho}=Y^{r_{j_0}}_{n_{j_0}}(f)_{p,\rho}$,


следовательно, выбрав $n$ так, что 


$\frac{1}{n+1}<\delta_{j_0}\le\frac{1}{n}$, 


по теореме 4 получим 


$$ 


\Omega^{r_{j_0}}_{i_{j_0}}(f)_{p,\rho}\preceq 


n^{-r_{j_0}}\sum^n_{s=1}s^{r_{j_0}-1} 


s^{-\alpha_{j_0}}\preceq\delta^{\alpha_{j_0}}_{j_0}.\quad\square


$$ 





Следствие 2 доказывается аналогично. 


 


Автор выражает благодарность В.М.\,Ф\"едорову 


за постановку задачи и внимание к работе. 


 


\begin{thebibliography}{9} 


 


\bibitem{1} 


Ф\"едоров В.М. {\em Приближение алгебраическими многочленами с весом 


Чебышева--Эрмита} 


// Изв. вузов. Математика. -- 1984. -- \No\,6. -- C.\,55--63. 


\bibitem{2} 


Алексеев Д.В. {\em Приближение полиномами функций одной переменной в 


метрике Чебышева--Эрмита} // Вестн. Моск. ун-та. -- 1997. -- 


\No\,6. -- C.\,68--71. 


\bibitem{3} 


Ржавинская Е.В. {\em О приближении функций двух переменных 


двумерным углом классических 


ортогональных полиномов}. -- Моск. ин-т электрон. техн. -- 


М., 1980. -- 42 с. -- Деп. в ВИНИТИ 30.09.80, \No\,4248-80. 


\bibitem{4} 


Стечкин С.Б., Субботин Ю.Н. {\em Сплайны в вычислительной 


математике}. -- М.: Наука, 1976. -- 248 с. 


 


\end{thebibliography} 


 


\vskip 0.5cm 


 


\post{\it Московский государственный}{\it Поступила} 


\post{\it университет}{29.10.1998} 


 


\label{end} 


\end{document} 


